\chapter{Lie Algebra, Lie Groups, and Compositional Maps}

Lie methods provide a useful change of viewpoint.  Instead of describing a
beam-line element only by the map that sends its entrance coordinates to its
exit coordinates, one describes an infinitesimal transformation---a
\emph{generator}---and obtains the finite map by exponentiation.  Linear
transfer matrices are the simplest example of this idea.  The same language
also applies to nonlinear Hamiltonian maps, with the matrix commutator
replaced by the Poisson bracket.



\section{Lie Algebra}

\subsection{Vector spaces and Lie algebras}

A vector space $V$ over a field $\mathbb{F}$ (here typically
$\mathbb{R}$ or $\mathbb{C}$) is a set whose elements can be added and
multiplied by scalars.  The result must remain in $V$, and these operations
obey the familiar linear rules: addition is associative and commutative,
there is a zero and an additive inverse, and scalar multiplication is
associative and distributes over both vector addition and scalar addition.
Consequently, if $X,Y\in V$ and $a,b\in\mathbb{F}$, then
$aX + bY\in V$.  This last closure test is often the most useful quick check in
practice.

A \emph{Lie algebra} $\mathfrak{g}$ over $\mathbb{F}$ is a vector space
equipped with a binary operation
\begin{equation}
    [\ ,\ ]:\mathfrak{g}\times\mathfrak{g}\longrightarrow\mathfrak{g}\,,
    \qquad (X,Y)\longmapsto[X,Y]\,,
\end{equation}
called the \emph{Lie bracket}.  For $X,Y,Z\in\mathfrak{g}$ and
$a,b\in\mathbb{F}$ it satisfies
\begin{align}
    [aX + bY,Z] & = a[X,Z] + b[Y,Z]\,,\notag\\
    [Z,aX + bY] & = a[Z,X] + b[Z,Y]
    &&\text{(bilinearity)},\label{eq:lie-bilinear}\\
    [X,X] & = 0
    &&\text{(alternation)},\label{eq:lie-alternation}\\
    [X,[Y,Z]] + [Y,[Z,X]] + [Z,[X,Y]] & = 0
    &&\text{(Jacobi identity)}\,.\label{eq:lie-jacobi}
\end{align}
Over the fields used here, $\mathbb R$ and $\mathbb C$, alternation is equivalent
to antisymmetry, $[X,Y] = -[Y,X]$.  Writing the axiom as alternation is also
correct in characteristic $2$, where a minus sign does not distinguish
antisymmetry from symmetry.
The bracket therefore measures non-commutativity.  If every bracket
vanishes, the algebra is called abelian.  Notice that a Lie algebra has two
quite different operations: vector addition and the bracket.  The bracket is
not in general an associative multiplication.


\subsection{Lie groups and representations}
A \emph{Lie group} $G$ is both a group and a smooth manifold, with smooth
group multiplication and inversion.  Its Lie algebra is the tangent space at
the identity,
\begin{equation}
    \mathfrak{g} = T_eG\,.
\end{equation}
A tangent vector $X\in T_eG$ is the initial velocity of a smooth curve
$g(t)$ through the identity: $g(0) = e$ and $\dot g(0) = X$.  Here $t$ is
simply a real curve parameter, defined at least on an open interval around
$t = 0$; it need not represent physical time.  Different curves with the same
initial velocity represent the same tangent vector.  Translating this vector
from the identity to every other group element produces a left-invariant vector
field. The usual commutator of these vector fields, evaluated back at $e$,
defines the Lie bracket.  Thus the algebra contains the infinitesimal
transformations of the group.  It determines the group locally near the
identity, but does not by itself contain all global topological information.

The group itself is usually \emph{not} a vector space.  Rotation matrices,
for example, form a group under matrix multiplication, but the sum of two
rotation matrices need not be a rotation matrix, and multiplying a rotation
matrix by zero leaves the group.  This is precisely why it is useful to move
to the tangent space: $T_eG$ is a vector space even when $G$ is curved.

If $\mathfrak{g}$ is finite-dimensional, choose a basis
$t_1,\ldots,t_n$.  These basis elements are commonly called the
\emph{generators}.  Every infinitesimal transformation is
\begin{equation}
    X \in \mathfrak{g} \;\Leftrightarrow\; X = x^at_a\,,
\end{equation}
where repeated indices are summed, and closure of the bracket gives
\begin{equation}
    [t_a,t_b] = f_{ab}^{\phantom{ab}c}t_c\,.
\end{equation}
The numbers $f_{ab}^{\phantom{ab}c}$ are the structure constants in this
basis.  Antisymmetry and the Jacobi identity impose corresponding relations
on them.  A different basis changes the structure constants but not the Lie
algebra.  The dimension of $\mathfrak{g}$ equals the dimension of $G$ as a
manifold.

A group or algebra can act on many different spaces.  A
\emph{representation} of a group on a vector space $V$ is a homomorphism
\footnote{
    The general linear group $\operatorname{GL}_n(\mathbb{F})$ is the set of
    $n\times n$ invertible matrices over the field $\mathbb{F}$, together
    with the operation of ordinary matrix multiplication.
}
\begin{equation}
    \rho:G\longrightarrow \operatorname{GL}(V)\,,
    \qquad \rho(g_1g_2) = \rho(g_1)\rho(g_2)\,,
\end{equation}
and a representation of its algebra is a linear map
\footnote{The result on the right is often written as $[\rho_*(X),\rho_*(Y)]$
using the ordinary commutator. One has hence to be careful to distinguish
it from the Lie bracket.}
\begin{equation}
    \rho_*:\mathfrak{g}\longrightarrow\mathfrak{gl}(V)\,,
    \qquad
    \rho_*([X,Y]) = \rho_*(X)\rho_*(Y) - \rho_*(Y)\rho_*(X)\,.
\end{equation}
A representation is
therefore a concrete realisation of abstract generators as matrices or
linear operators.  It is \emph{faithful} if different algebra elements give
different operators.  Much of Lie theory in accelerator physics consists of
choosing the representation best suited to the calculation: matrices for
linear optics, differential operators for nonlinear maps, or finite
truncations acting on Taylor coefficients (TPSAs) for numerical
implementations of nonlinear maps.


\subsection{Example: real \texorpdfstring{$2\times2$}{2 x 2} matrices}

Consider
\begin{equation}
    G = \operatorname{GL}_2(\mathbb{R})
      = \{A\in\mathbb{R}^{2\times2}:\det A\ne0\}\,.
\end{equation}
This is a Lie group under matrix multiplication.  It is not a vector space:
the zero matrix is excluded, and two invertible matrices can have a singular
sum.  A curve through the identity can be written
$A(t) = I + tX + \mathcal{O}(t^2)$, where $X$ is an arbitrary real $2\times2$
(potentially singular) matrix.  Hence
\begin{equation}
    \mathfrak{gl}(2,\mathbb{R}) = \mathbb{R}^{2\times2}\,,
    \qquad [X,Y] = XY - YX\,.
\end{equation}
The commutator obeys Eqs.~\eqref{eq:lie-bilinear}--\eqref{eq:lie-jacobi}.
Using the four matrices $E_{ij}$, with a one in row $i$, column $j$, and
zeros elsewhere, as generators gives
\begin{equation}
    [E_{ij},E_{kl}] = \delta_{jk}E_{il}-\delta_{li}E_{kj}\,.
\end{equation}
For example,
$[E_{12},E_{21}] = E_{11}-E_{22}$.  This is the elementary model for the
same commutator structure encountered with larger transfer matrices.


\subsection{Example: \texorpdfstring{$\mathfrak{su}(N)$}{su(N)}}

The special unitary group is
\begin{equation}
    \operatorname{SU}(N) = \{U\in\mathbb{C}^{N\times N}:U^\dagger U
                         = I,\ \det U = 1\}\,.
\end{equation}
Put $U(t) = I + tX + \mathcal{O}(t^2)$ and differentiate both defining
conditions at $t = 0$.  One obtains
\begin{equation}
    \mathfrak{su}(N) = \{X\in\mathbb{C}^{N\times N}:X^\dagger
                     = -X,\ \Tr X = 0\}\,,
    \qquad [X,Y] = XY - YX\,.
\end{equation}
It is a real vector space of dimension $N^2-1$.  Physicists often prefer
Hermitian generators $t_a$ and write $X_a = -\ii t_a$.  With the conventional
normalisation,
\begin{equation}
    [t_a,t_b] = \ii f_{ab}^{\phantom{ab}c}t_c\,,
    \qquad
    [X_a,X_b] = f_{ab}^{\phantom{ab}c}X_c\,.
\end{equation}
For $\operatorname{SU}(2)$ one may take $t_a = \sigma_a/2$, where the
$\sigma_a$ are the Pauli matrices, and
$f_{ab}^{\phantom{ab}c} = \epsilon_{abc}$.  This example is useful because it
separates the abstract real Lie algebra from two common matrix conventions:
anti-Hermitian generators are closed under the ordinary commutator, whereas
Hermitian generators introduce an explicit factor of $\ii$.



\section{Hamiltonian Maps in Six-Dimensional Phase Space}

The application of Lie algebra to nonlinear accelerator physics takes us
one step beyond ordinary matrix spaces, and requires a bit of differential
geometry to lay the necessary mathematical foundation. One can skip some of
these parts and jump directly to the relevant part on the Lie algebra of
Hamiltonian maps in Sec.~\ref{sec:hamiltonian-lie-algebra}.


\subsection{Manifolds}
A smooth $n$-dimensional \emph{manifold} $M$ is a space that looks locally
like $\mathbb{R}^n$: every point has a neighbourhood that can be described
by $n$ coordinates, and changes between overlapping coordinate systems are
smooth.  Coordinates therefore let us perform calculations locally, but the
coordinate tuple assigned to a point should not be confused with the point
itself.

For a particle in three degrees of freedom, the familiar canonical phase
space is
\begin{equation}
    M = \mathbb{R}^6,\qquad
    \vec z = (q_1,p_1,q_2,p_2,q_3,p_3)^T\,.
\end{equation}
This particular manifold \emph{is} a vector space: componentwise addition
and scalar multiplication remain in $\mathbb{R}^6$, and all vector-space
axioms follow component by component from those of $\mathbb{R}$.  A general
manifold, however, need not be a vector space.  For instance, the
configuration space of a plane pendulum is the circle $S^1$, and its phase
space is
\begin{equation}
    T^*S^1 \simeq S^1\times\mathbb{R}\,,
\end{equation}
which is a cylinder: the angle is periodic, whereas its conjugate momentum
is real.  By itself, the manifold structure supplies neither addition nor
scalar multiplication of points. We could define addition as, for instance,
\begin{equation}
    \left(\theta_1, p_1\right) + \left(\theta_2, p_2\right)
    = \left(\theta_1 + \theta_2 \;(\text{mod } 2\pi), p_1 + p_2\right)\,,
\end{equation}
which works beautifully. But notice something subtle: here we have implicitly
chosen which point of the circle corresponds to $\theta = 0$ to deal with the
periodicity of the angle. And this is the key point: the manifold structure by
itself supplies no preferred addition law for points.  A coordinate-independent
addition law can certainly exist, as it does when the manifold is also a Lie
group, but it is \emph{additional structure} rather than part of the definition
of a manifold. More fundamentally, every finite-dimensional real vector space
is topologically equivalent to some $\mathbb{R}^n$ and can be continuously
contracted to a point, whereas a loop around the cylinder cannot be contracted.
Thus, being locally Euclidean does not make a manifold a vector space.


\subsection{Tangent space}
Although points of a general manifold cannot be added, the infinitesimal
directions at any fixed point can.  These directions form the \emph{tangent
space} $T_p M$ at $p\in M$.  One way to define a tangent vector at $p$ is by
a smooth curve $\gamma(t)$ through that point, with $\gamma(0) = p$: two curves
represent the same tangent vector if their derivatives at $t = 0$ agree in one
(and hence every) local coordinate system.  In local coordinates
$(x^1,\ldots,x^n)$, a tangent vector can be written
\begin{equation}
    v = \sum_{i = 1}^n v^i\left.\frac{\partial}{\partial x^i}\right|_p\,.
\end{equation}
Its components depend on the chosen coordinates, but the vector itself does
not.

Crucially, $T_p M$ \emph{is} an $n$-dimensional vector space, even when $M$
is not: tangent vectors based at the same point may be added and multiplied
by scalars.  The collection of tangent spaces at all points is the
\emph{tangent bundle}
\begin{align}
    T M & = \bigsqcup_{p\in M} T_p M\,, \\
    & = \{(p, v) \mid p \in M, v \in T_p M\}\,.
\end{align}
The \emph{cotangent space} $T_p^* M$ is the dual vector space of $T_p M$; its
elements are linear maps $T_p M\to\mathbb{R}$.  In the coordinates above it
has the dual basis $(\dee x^1,\ldots,\dee x^n)$.  Collecting these dual spaces
gives the \emph{cotangent bundle}
\begin{align}
    T^* M & = \bigsqcup_{p\in M} T_p^* M\,, \\
    & = \{(p, v^*) \mid p \in M, v^* \in T_p^* M\}\,.
\end{align}
This explains the notation used for the pendulum: a point of $T^*S^1$ is a
pair $(\theta,p_\theta)$, where $\theta\in S^1$ specifies the configuration
and $p_\theta\,\dee\theta\in T_\theta^*S^1$ is its conjugate momentum.
\paragraph{Practical interpretation in phase space.}
For accelerator calculations, these abstract spaces have a very concrete
meaning.  If $z = (q_1,p_1,\ldots,q_n,p_n)\in M$ is a phase-space point, then a
tangent vector $v\in T_zM$ is represented in canonical coordinates by an
infinitesimal phase-space variation
\begin{equation}
    \delta\vec z = (\delta q_1,\delta p_1,\ldots,
    \delta q_n,\delta p_n)^T\,.
\end{equation}
Equivalently, it may be the velocity $\dot{\vec z}(0)$ of a curve through
$z$.  Thus the tangent vectors of phase space are \emph{not} the canonical
momenta: the $p_i$ are already coordinates of the point $z\in M$, whereas a
tangent vector describes how the complete phase-space point changes.  For
example, the Jacobian of a tracking map acts precisely on such variations,
$\delta\vec z_{\rm out} = D\mathcal M_z\,\delta\vec z_{\rm in}$.

A cotangent vector $\alpha\in T_z^*M$ is instead a linear rule that assigns a
number to each infinitesimal variation.  In coordinates,
\begin{equation}
    \alpha = \alpha_i\,\dee z^i,
    \qquad
    \alpha(v) = \alpha_i v^i\,.
\end{equation}
The differential of a scalar function is the most common example:
$\dee f = (\partial f/\partial z^i)\dee z^i$ and
$\dee f(v) = v^i\partial_i f$, the directional derivative of $f$ along $v$.
Writing a covector as a row vector or ``gradient'' is therefore a coordinate
representation; identifying it with a vector by a dot product additionally
uses a metric and should not be confused with the intrinsic vector--covector
duality.

There is one related use of the word momentum that is important but lives one
level lower.  If the phase space itself is a cotangent bundle $M = T^*Q$, then
a phase-space point is $(q,p)$ with $q\in Q$ and $p\in T_q^*Q$.  In that sense
the canonical momenta are the components of a covector on the
\emph{configuration space} $Q$; they are still not tangent vectors of the
phase space $M$.


\subsection{Symplectic manifolds}
Phase space carries yet another layer of structure.  At a point $p\in M$, an
alternating covariant 2-tensor is an element
\begin{equation}
    \omega_p\in\bigwedge^2 T_p^*M\,.
\end{equation}
Here $\bigwedge^2 T_p^*M$ is the vector space of antisymmetric bilinear
functions
\begin{equation}
    \omega_p:T_pM\times T_pM\longrightarrow\mathbb{R},
    \qquad \omega_p(u,v) = -\omega_p(v,u)\,.
\end{equation}
A \emph{differential 2-form} $\omega$ on $M$ is a choice of such a tensor
$\omega_p$ at every $p$ that varies smoothly with $p$.  The vector space of
all differential 2-forms on $M$ is denoted by $\Omega^2(M)$.  The pair
$(M,\omega)$ is a \emph{symplectic manifold} if $\omega$ is closed and
non-degenerate:
\begin{align}
    \dee\omega & = 0\,, \\
    \omega_p(v,w) = 0\ \forall w\in T_pM
    &\quad\Longleftrightarrow\quad v = 0\,.
\end{align}
Here $\dee\omega$ is the exterior derivative of $\omega$, and the second
condition is required at every $p\in M$.
An antisymmetric bilinear form can be non-degenerate only in even dimensions.
Darboux's theorem guarantees that, near every point of a
$2n$-dimensional symplectic manifold, there are canonical coordinates in
which
\begin{equation}
    \omega = \sum_{j = 1}^{n}\dee q_j\wedge\dee p_j\,.
\end{equation}
For the pendulum this is $\omega = \dee\theta\wedge\dee p_\theta$; the same
local formula works even though $\theta$ is not a global real coordinate.
For the special case $M = \mathbb R^6$, there is a natural identification
\begin{equation}
    T_z\mathbb R^6\simeq\mathbb R^6\,.
\end{equation}
at every $z\in \mathbb R$.  Therefore, we can write
\begin{equation}
    \omega:\mathbb R^6\times\mathbb R^6\to\mathbb R, \quad
    \omega(u,v) = u^T S v\,,
\end{equation}
with
\begin{equation}
    S = \begin{pmatrix}
    0 & 1 & 0 & 0 & 0 & 0\\
    -1 & 0 & 0 & 0 & 0 & 0\\
    0 & 0 & 0 & 1 & 0 & 0\\
    0 & 0 & -1 & 0 & 0 & 0\\
    0 & 0 & 0 & 0 & 0 & 1\\
    0 & 0 & 0 & 0 & -1 & 0
    \end{pmatrix}\,.
\end{equation}
We can hence think of $S$ as the coordinate representation of $\omega$.


\subsection{Symplectic maps}
A canonical, or symplectic, map $\mathcal{M}:M\to M$ is a diffeomorphism
that preserves the symplectic form.  The easiest way to understand this is to
follow two infinitesimal phase-space variations.  At $z\in M$, the
differential (or \emph{pushforward})
\begin{equation}\label{eq:pushforward}
    D\mathcal{M}_z:T_zM\longrightarrow T_{\mathcal{M}(z)}M\,.
\end{equation}
propagates a tangent vector $u$ at the entrance into the tangent vector
$D\mathcal M_z u$ at the image point.  In accelerator language this is the
linearised tracking map about the reference point $z$.

The symplectic form at the image point can therefore be evaluated on the two
transported vectors,
\begin{equation}
    \omega_{\mathcal M(z)}
    \bigl(D\mathcal M_z u,D\mathcal M_z v\bigr)\,.
\end{equation}
This quantity depends on the original point $z$ and on the original vectors
$u,v\in T_zM$, so it defines a new 2-form at $z$.  That new form is, by
definition, the \emph{pullback} of $\omega$ by $\mathcal M$:
\begin{equation}\label{eq:pullback-two-form}
    (\mathcal M^*\omega)_z(u,v)
    \equiv
    \omega_{\mathcal M(z)}
    \bigl(D\mathcal M_z u,D\mathcal M_z v\bigr)\,.
\end{equation}
Thus $D\mathcal M$ pushes vectors forward, while $\mathcal M^*$ pulls forms
back so that they can again be compared at the original point.  The map is
symplectic precisely when this comparison shows no change,
\begin{equation}\label{eq:symplectic-condition}
    \mathcal M^*\omega = \omega\,,
\end{equation}
or, point by point,
\begin{equation}
    \omega_{\mathcal M(z)}
    \bigl(D\mathcal M_z u,D\mathcal M_z v\bigr) = \omega_z(u,v)\,,
\end{equation}
for every $z$ and every $u,v\in T_zM$.  Nothing physically moves ``backwards''
in a pullback: it is simply the bookkeeping operation that expresses a form
from the output point in terms of input tangent vectors.

In canonical coordinates on the $2n$-dimensional phase space, the
differential $D\mathcal M_{\vec z}$ is represented by the usual Jacobian
matrix, with entries
\begin{equation}
    \bigl[D\mathcal M_{\vec z}\bigr]_{ij}
    = \frac{\partial \mathcal M_i}{\partial z_j}\,.
\end{equation}
Thus, if $\vec u$ is the coordinate vector representing $u\in T_zM$, then
$D\mathcal M_{\vec z}\,\vec u$ represents
$D\mathcal M_z u\in T_{\mathcal M(z)}M$. Since the matrix of $\omega$ is
the constant matrix $S$ in these coordinates, the preceding identity becomes,
for arbitrary coordinate vectors $\vec u$ and $\vec v$,
\begin{equation}
    \bigl(D\mathcal M_{\vec z}\,\vec u\bigr)^T S
    \bigl(D\mathcal M_{\vec z}\,\vec v\bigr)
    = \vec u^T S\vec v\,.
\end{equation}
Equivalently,
\begin{equation}
    \vec u^T\left[
        \bigl(D\mathcal M_{\vec z}\bigr)^T S D\mathcal M_{\vec z} - S
    \right]\vec v = 0\,,
\end{equation}
for all $\vec u$ and $\vec v$.  The matrix in brackets must therefore
vanish, which gives the familiar Jacobian condition
\begin{equation}\label{eq:symplectic-jacobian}
    \bigl(D\mathcal M_{\vec z}\bigr)^T S D\mathcal M_{\vec z} = S\,.
\end{equation}
Consequently, in canonical coordinates,
Eq.~\eqref{eq:symplectic-jacobian} is equivalent to the geometric
symplecticity condition~\eqref{eq:symplectic-condition}.

For a \emph{linear} map $\vec z\mapsto R\vec z$, the Jacobian is simply the
constant matrix $R$.  The linear symplectic maps therefore form the
\emph{symplectic group}
\begin{equation}
    \operatorname{Sp}(2n,\mathbb R) = \{R\in\operatorname{GL}(2n,\mathbb R)
    \mid R^T S R = S\}\,.
\end{equation}
This is the group containing the familiar linear transfer matrices of
Hamiltonian beam optics.  To find its Lie algebra, write a curve through the
identity as $R(t) = I + tA + \mathcal O(t^2)$ and differentiate
$R(t)^TSR(t) = S$ at $t = 0$.  One obtains
\begin{equation}\label{eq:symplectic-lie-algebra}
    \mathfrak{sp}(2n,\mathbb R) = \{A\in\mathbb R^{2n\times2n}
    \mid A^TS + SA = 0\}\,.
\end{equation}
Equivalently, $SA$ is symmetric.  This provides the direct linear analogue of
the nonlinear construction below: a symplectic matrix is a finite canonical
transformation, while an element of $\mathfrak{sp}(2n,\mathbb R)$ is its
infinitesimal generator.  For quadratic Hamiltonians the Hamiltonian vector
field is linear, so exponentiating that vector field reproduces precisely a
matrix in $\operatorname{Sp}(2n,\mathbb R)$.  Higher-order Hamiltonians are
the nonlinear extension of the same idea.


\subsection{Hamiltonian Lie algebra and Poisson algebra}
\label{sec:hamiltonian-lie-algebra}
The infinitesimal generators of Hamiltonian maps are smooth functions on
phase space.  Let $C^\infty(M)$ denote all smooth real-valued functions on
$M$.  It is a vector space under pointwise operations,
\begin{equation}
    (af + bg)(z) = af(z) + bg(z)\,.
\end{equation}
Indeed, a linear combination of smooth functions is smooth, so the set is
closed; all the remaining vector-space laws follow from the corresponding
laws for real numbers at every point $z$.  This is also an
infinite-dimensional vector space.  For example, on a canonical coordinate
patch the functions $1,q_1,q_1^2,\ldots$ (multiplied by a smooth cutoff if
the patch is not global) contain arbitrarily large linearly independent
subsets.

The relevant functions here are \emph{smooth Hamiltonians}, not necessarily
``integrable functions.''  An integrable Hamiltonian in the Liouville sense
is a Hamiltonian system with sufficiently many independent invariants in
involution; such Hamiltonians do not form the vector space used below.
Likewise, mere integrability in the measure-theoretic sense would not provide
the derivatives needed for a Poisson bracket.

In canonical coordinates the Poisson bracket is
\begin{equation}\label{eq:canonical-poisson-bracket}
    \poissonfull{f}{g} = \sum_{j = 1}^{n}\left(
    \partder{f}{q_j}\partder{g}{p_j}
    -\partder{f}{p_j}\partder{g}{q_j}
    \right)\,.
\end{equation}
This is the convention used by Dragt and Forest: the first argument is
differentiated with respect to the coordinate $q_j$ in the first term.  In
particular,
\begin{equation}
    \poisson{q_i}{p_j} = \delta_{ij},\qquad
    \poisson{q_i}{q_j} = \poisson{p_i}{p_j} = 0\,.
\end{equation}
It is bilinear and antisymmetric, satisfies the Jacobi identity, and maps
two smooth functions to another smooth function.  Therefore
$(C^\infty(M),\poissonfull{\slot}{\slot})$ is a Lie algebra.  A constant has
zero Poisson bracket with every function and generates no motion, so
Hamiltonians that differ only by a constant describe the same infinitesimal
map.  On a connected phase space the non-redundant algebra is consequently
\begin{equation}\label{eq:hamiltonian-lie-algebra}
    \mathfrak{ham}(M)\simeq C^\infty(M)/\mathbb{R}\,.
\end{equation}
This quotient is again a vector space: define
$\lfloor f\rfloor + \lfloor g\rfloor = \lfloor f + g\rfloor$ and
$a\lfloor f\rfloor = \lfloor a f\rfloor$, where $\lfloor f\rfloor$ denotes
all functions that differ from $f$ by a constant.  Changing representatives
only changes the result by a constant, so these operations are well-defined.
The same argument and $\poisson{f + c}{g + d} = \poisson{f}{g}$ show that
the Poisson bracket is well-defined on the quotient.  The algebra remains
infinite-dimensional.  With the conventions used here, the bracket-preserving
identification with the tangent vector fields is
\begin{equation}
    \lfloor f\rfloor \longmapsto -X_f\,,
    \qquad [-X_f,-X_g] = -X_{\poisson{f}{g}}\,.
\end{equation}
One may instead identify $\lfloor f\rfloor$ with $X_f$.  The corresponding
function-space bracket then acquires an overall minus sign.  This harmless
choice is the origin of several sign conventions in the Lie-map literature.

There is a small global qualification behind
Eq.~\eqref{eq:hamiltonian-lie-algebra}. The tangent algebra to all symplectic
maps consists of symplectic vector fields; a symplectic vector field need only
be locally Hamiltonian.  On the usual $\mathbb{R}^6$ phase space every such
field of interest is generated by a global Hamiltonian, while on a manifold
with non-trivial topology the Hamiltonian maps can form a proper subgroup of
all symplectic maps.  For local accelerator maps, $C^\infty(M)/\mathbb{R}$ is
the algebra normally used.

The full function space has still more structure.  With ordinary pointwise
multiplication it is a commutative associative algebra, and the bracket
satisfies the Leibniz rule
\begin{equation}\label{eq:poisson-leibniz}
    \poisson{f}{gh} = \poisson{f}{g}h + g\poisson{f}{h}\,.
\end{equation}
A commutative algebra equipped with a Lie bracket satisfying
Eq.~\eqref{eq:poisson-leibniz} is called a \emph{Poisson algebra}.  Thus
$C^\infty(M)$ is both a commutative algebra and a Lie algebra, with the two
products tied together by the Leibniz rule.  Strictly speaking, the quotient
$C^\infty(M)/\mathbb{R}$ is a Lie algebra but not a Poisson algebra under
pointwise multiplication, because constants are not an ideal for that
multiplication.  In calculations one therefore keeps representatives in
$C^\infty(M)$ and discards irrelevant constant terms when interpreting the
generator.

This Poisson-algebra structure is particularly useful for accelerator
physics.  It permits brackets of polynomial Hamiltonians to be evaluated by
ordinary differentiation and multiplication, it extends immediately to
Taylor and truncated-power-series representations, and the Leibniz rule
controls the action of a generator on products of observables.  At the same
time, working with a Hamiltonian generator rather than arbitrary Taylor
coefficients keeps the underlying map canonical.



\section{Exponentiation}

\subsection{From an infinitesimal generator to a finite transformation}

For a Lie group $G$ with algebra $\mathfrak{g}$, the exponential map
\begin{equation}
    \exp:\mathfrak{g}\longrightarrow G\,.
\end{equation}
is defined through one-parameter subgroups.  For every $X\in\mathfrak g$
there is a unique one-parameter subgroup $\gamma_X:\mathbb R\to G$ with
\begin{equation}
    \gamma_X(0) = e,\qquad
    \dot\gamma_X(0) = X,\qquad
    \gamma_X(t + s) = \gamma_X(t)\gamma_X(s)\,.
\end{equation}
The parameter is $t\in\mathbb R$, not merely $t\ge0$: negative values give
the inverse transformation,
$\gamma_X(-t) = \gamma_X(t)^{-1}$.  At this abstract stage $t$ is only a group
parameter; depending on the application it may later be physical time,
longitudinal position, an angle, a magnet strength multiplied by a length, or
simply a dimensionless interpolation parameter.  One defines
\begin{equation}
    \exp(tX)\equiv\gamma_X(t),\qquad \exp(X) = \gamma_X(1)\,.
\end{equation}
Thus exponentiation means following the same infinitesimal generator $X$ for
a finite amount of its parameter.  This geometric construction is the
definition; a power series appears only after choosing a representation in
which the generator is a linear operator.

For a matrix Lie group, $X$ is a matrix and the group exponential is exactly
the convergent matrix series
\begin{equation}\label{eq:matrix-exponential}
    \exp X = I + X + \frac{X^2}{2!} + \frac{X^3}{3!} + \cdots\,.
\end{equation}
More generally, let
\begin{equation}
    \rho:G\to\operatorname{GL}(V)\,.
\end{equation}
be a representation of the group.  Its differential at the identity,
\begin{equation}
    \rho_*\equiv(\dee\rho)_e:
    \mathfrak g\to\mathfrak{gl}(V)\,,
\end{equation}
is the corresponding Lie-algebra representation introduced above.  It maps
the abstract infinitesimal generator $X$ to the linear operator
$\rho_*(X)$ acting on $V$.  Group exponentiation and operator exponentiation
then commute with the representation:
\begin{equation}\label{eq:representation-exponential}
    \rho(\exp tX) = \exp\!\bigl(t\,\rho_*(X)\bigr)\,.
\end{equation}
This equation is not an additional definition: it states that representing a
finite group transformation gives the same result as first representing its
generator and then exponentiating that operator.

For bounded operators the analogous operator series converges.  For
unbounded differential operators, or for formal Taylor maps, its domain and
convergence require care and the series is often used order by order as a
formal series.  Moreover, the exponential need not be one-to-one or cover
the entire group, although for a finite-dimensional Lie group its image
contains a neighbourhood of the identity.


\subsection{Example: \texorpdfstring{$\operatorname{SU}(N)$}{SU(N)}}

If $X\in\mathfrak{su}(N)$, Eq.~\eqref{eq:matrix-exponential} gives an
element of $\operatorname{SU}(N)$.  Indeed,
\begin{equation}
    (\e^X)^\dagger\e^X = \e^{X^\dagger}\e^X = \e^{-X}\e^X = I\,,
    \qquad
    \det(\e^X) = \e^{\Tr X} = 1\,.
\end{equation}
With Hermitian physics generators this is commonly written
\begin{equation}
    U(\vec\theta) = \exp(-\ii\theta^aT_a)\,.
\end{equation}
The generator is linear in the parameters, while the exponential builds the
finite, generally non-commuting transformation.


\subsection{Example: Hamiltonian flow}

The preceding abstract construction becomes especially concrete in phase
space.  Let $f\in C^\infty(M)$.  Its Hamiltonian vector field $X_f$ is the
vector field defined, for the symplectic convention used here, by
\begin{equation}\label{eq:hamiltonian-vector-field-geometric}
    \iota_{X_f}\omega = \dee f\,,
\end{equation}
where $\iota_{X_f}$ denotes insertion of $X_f$ into the first argument of the
2-form.  In canonical coordinates this gives
\begin{equation}\label{eq:hamiltonian-vector-field}
    X_f = \sum_{j = 1}^{n}\left(
    \partder{f}{p_j}\partder{}{q_j}
    -\partder{f}{q_j}\partder{}{p_j}
    \right),
    \qquad
    X_f g = \poisson{g}{f} = -\poisson{f}{g}\,.
\end{equation}
Thus $X_f(z)\in T_zM$ gives an infinitesimal direction of motion at every
phase-space point.

The \emph{flow} of this vector field is the family of point maps
\begin{equation}
    \Phi_f^t:M\to M\,.
\end{equation}
defined by the initial-value problem
\begin{equation}\label{eq:hamiltonian-flow-definition}
    \totder{}{t}\Phi_f^t(z) = X_f\!\left(\Phi_f^t(z)\right),
    \qquad
    \Phi_f^0(z) = z\,.
\end{equation}
For each initial point $z$, the curve $t\mapsto\Phi_f^t(z)$ is the integral
curve of $X_f$ through $z$.  If $f = H$ is the physical Hamiltonian and $t$ is
physical time, Eq.~\eqref{eq:hamiltonian-flow-definition} is simply
Hamilton's equations,
\begin{equation}
    \dot q_j = \partder{H}{p_j},\qquad
    \dot p_j = -\partder{H}{q_j}\,,
\end{equation}
or, for any observable $g$ with no explicit $t$ dependence,
\begin{equation}\label{eq:observable-hamilton-equation}
    \totder{g}{t} = X_Hg = \poisson{g}{H}\,.
\end{equation}
In accelerator applications the evolution parameter need not be time; after
a suitable Hamiltonian formulation it is often the longitudinal coordinate
$s$.  Mathematically nothing changes.

For an autonomous generator, the flow satisfies
\begin{equation}
    \Phi_f^{t + s} = \Phi_f^t\circ\Phi_f^s,
    \qquad
    (\Phi_f^t)^{-1} = \Phi_f^{-t}\,,
\end{equation}
whenever the relevant solutions exist.  It is therefore a one-parameter
(local, or global when the vector field is complete) group of
diffeomorphisms, and the group exponential is its unit-parameter map,
\begin{equation}
    \exp(tX_f) = \Phi_f^t,\qquad \exp(X_f) = \Phi_f^1\,.
\end{equation}
This is the precise meaning of exponentiating a Hamiltonian vector field.

The flow is symplectic.  Cartan's formula and
Eq.~\eqref{eq:hamiltonian-vector-field-geometric} give
\begin{equation}
    \mathcal L_{X_f}\omega = \dee(\iota_{X_f}\omega) + \iota_{X_f}(\dee\omega)
                           = \dee^2f = 0\,.
\end{equation}
Using the standard identity for differentiation of a pulled-back form along
a flow,
\begin{equation}
    \totder{}{t}(\Phi_f^t)^*\omega = (\Phi_f^t)^*(\mathcal L_{X_f}\omega) = 0\,,
\end{equation}
so $(\Phi_f^t)^*\omega = \omega$ because $\Phi_f^0$ is the identity.  In
words: infinitesimal Hamiltonian motion preserves the symplectic form, hence
so does the finite map obtained by following that motion.

If the generator depends explicitly on the evolution parameter,
$f = f_t$, the maps from different parameter values are generated by different
vector fields and do not in general form $\exp(tX_f)$ for one fixed $X_f$.
The solution is then a parameter-ordered exponential (or, equivalently, the
solution of the non-autonomous flow equation).  This is analogous to a
sequence of non-commuting beam-line elements, which cannot generally be
replaced by the exponential of their simple sum.


\subsection{The adjoint representation}

Every Lie group has a natural representation on its own algebra.  Conjugate
a group element close to the identity by $g\in G$ and differentiate at the
identity:
\begin{equation}
    \operatorname{Ad}_gX = \left.\totder{}{t}\right|_{t = 0}
    g\exp(tX)g^{-1}\,.
\end{equation}
This is the \emph{adjoint representation} of the group.  For a matrix group,
\begin{equation}
    \operatorname{Ad}_gX = gXg^{-1}\,.
\end{equation}
Differentiating once more gives the adjoint representation of the algebra,
\begin{equation}\label{eq:adjoint-algebra}
    \operatorname{ad}_X:\mathfrak{g}\longrightarrow\mathfrak{g}\,,
    \qquad \operatorname{ad}_X(Y) = [X,Y]\,.
\end{equation}
The Jacobi identity is exactly the statement
\begin{equation}\label{eq:ad-commutator}
    [\operatorname{ad}_X,\operatorname{ad}_Y] = \operatorname{ad}_{[X,Y]}\,.
\end{equation}
Thus the adjoint operators form a Lie algebra under the operator commutator.
More precisely, $\operatorname{ad}_X$ is an element of
$\mathfrak{gl}(\mathfrak{g})$, not in general an element of the original
$\mathfrak{g}$; the image $\operatorname{ad}(\mathfrak{g})$ is a Lie
subalgebra of operators.  Elements in the centre of $\mathfrak{g}$ act as
zero, so the adjoint representation is faithful precisely when the centre is
trivial.

Exponentiating this representation gives
\begin{equation}\label{eq:exp-adjoint}
    \operatorname{Ad}_{\exp X} = \exp(\operatorname{ad}_X)
    = I + \operatorname{ad}_X + \frac{\operatorname{ad}_X^2}{2!} + \cdots\,.
\end{equation}
This is now an ordinary exponential series of linear operators acting on
algebra elements.  For $\operatorname{SU}(N)$,
\begin{equation}
    \operatorname{Ad}_U(Y) = UYU^\dagger\,,
    \qquad \operatorname{ad}_X(Y) = [X,Y]\,,
\end{equation}
and Eq.~\eqref{eq:exp-adjoint} states
$\e^XY\e^{-X} = \e^{\operatorname{ad}_X}Y$.
In the generator basis $X_a = -\ii T_a$, the matrices of the adjoint
representation have components
$(\operatorname{ad}_{X_a})_b^{\phantom{b}c} = f_{ab}^{\phantom{ab}c}$\,.


\subsection{Lie operators on phase-space functions}

The adjoint representation of the Hamiltonian algebra is the Lie operator
used in nonlinear beam dynamics:
\begin{equation}\label{eq:colon-definition}
    \Lie{f} \,g \equiv \operatorname{ad}_f(g) = \poisson{f}{g}\,.
\end{equation}
It is a linear differential operator on functions, and the Jacobi identity
gives
\begin{equation}
    [\Lie{f},\Lie{g}] = \Lie{\poisson{f}{g}}\,.
\end{equation}
Constants disappear automatically, so this is a representation of
$C^\infty(M)/\mathbb{R}$.  The exponentiated Lie operator is
\begin{equation}\label{eq:lie-exponential}
    \e^{\lie{f}} g = g + \poisson{f}{g}
    + \frac{1}{2!}\poisson{f}{\poisson{f}{g}} + \cdots\,.
\end{equation}
Unlike the abstract definition of the group exponential, this expression
\emph{is} an exponential power series: it is the operator exponential in the
adjoint representation.

The relation to the point flow is now immediate.  With the Dragt--Forest
convention of Eq.~\eqref{eq:canonical-poisson-bracket},
\begin{equation}
    \Lie{f}g = \poisson{f}{g} = -X_fg\,.
\end{equation}
Therefore the operator exponential acts on a function by pullback along the
\emph{inverse} Hamiltonian flow,
\begin{equation}\label{eq:lie-flow-relation}
    \e^{t\,\lie{f}}g = g\circ\Phi_f^{-t} = (\Phi_f^{-t})^*g\,.
\end{equation}
In particular, applying it componentwise to the coordinate functions gives
\begin{equation}
    \Phi_H^t(\vec z) = \e^{-t\,\lie{H}}\vec z\,.
\end{equation}
This minus sign is not a dynamical convention hidden in Hamilton's equations;
it comes solely from defining the Lie operator with the generator in the
\emph{first} Poisson-bracket slot, $\Lie{f}g = \poisson{f}{g}$, whereas the
Hamiltonian vector field acts as $X_fg = \poisson{g}{f}$.  This is the
Dragt--Forest convention used throughout this chapter.

Because of the Leibniz rule, $\Lie{f}$ is a derivation.  Its exponential is
therefore an algebra automorphism:
\begin{equation}
    \e^{\lie{f}}(gh) = (\e^{\lie{f}}g)(\e^{\lie{f}}h)\,,
    \qquad
    \poisson{\e^{\lie{f}}g}{\e^{\lie{f}}h} = \e^{\lie{f}}\poisson{g}{h}\,.
\end{equation}
Acting on the coordinate functions gives a canonical map,
\begin{equation}
    \mathcal{M}_f(\vec z) = \e^{\lie{f}}\vec z = \Phi_f^{-1}(\vec z)\,,
\end{equation}
where the exponential acts on every component of $\vec z$ and the final
equality refers to unit flow parameter.  Thus a Lie map with generator $f$
is the inverse unit-parameter flow of the Hamiltonian vector field $X_f$ under
the convention adopted here.  Products of
these exponentials are the \emph{compositional maps} used in Lie
factorisations~\cite{DragtFinn1976,DragtForest1983}.  Non-commutativity is
retained through the Baker--Campbell--Hausdorff series; formally,
\begin{equation}
    \e^{\lie{f}}\e^{\lie{g}} = \e^{\lie{h}},\qquad
    h = f + g + \frac{1}{2}\poisson{f}{g} + \cdots\,,
\end{equation}
up to an irrelevant additive constant in $h$.


\subsection{Why map order can appear to be reversed}

Two related objects are frequently denoted by the same symbol in the
accelerator literature:
\begin{enumerate}
 \item a \emph{point map} $\mathcal{A}:M\to M$, which takes a phase-space
       point as its argument; and
 \item the corresponding \emph{coordinate-substitution operator}
       $\widehat{\mathcal{A}}$, which acts on a function according to
       \begin{equation}
    \widehat{\mathcal{A}}g = g\circ\mathcal{A}\,.
\end{equation}
\end{enumerate}
Even if $\mathcal{A}$ is nonlinear, $\widehat{\mathcal{A}}$ is a linear
operator on the vector space of functions.  In particular,
$\widehat{\mathcal{M}_f} = \e^{\lie{f}}$.

Suppose a particle first passes through $\mathcal{A}$ and then through
$\mathcal{B}$.  For point maps with explicit arguments,
\begin{equation}\label{eq:point-map-order}
    \vec z_1 = \mathcal{A}(\vec z_0)\,,\qquad
    \vec z_2 = \mathcal{B}(\vec z_1)
             = \left(\mathcal{B}\circ\mathcal{A}\right)(\vec z_0)\,.
\end{equation}
The rightmost point map acts first.  Matrices multiplying a column vector
follow the same rule: $\vec z_2 = BA\vec z_0$.

Substitution operators compose in the opposite order:
\begin{equation}\label{eq:pullback-order}
    \widehat{\mathcal{B}\circ\mathcal{A}}\,g
    = g\circ\mathcal{B}\circ\mathcal{A}
    = \widehat{\mathcal{A}}\,\widehat{\mathcal{B}}\,g\,.
\end{equation}
For the coordinate functions themselves this becomes
\begin{equation}
    \vec z_2 = \left(\widehat{\mathcal{A}}\,
    \widehat{\mathcal{B}}\,\vec z\right)(\vec z_0)\,.
\end{equation}
The labels in the operator product therefore follow the physical beam-line
order from left to right, even though, as operators acting on the object to
their right, the rightmost operator is evaluated first.  This pullback
reversal is the source of the apparently conflicting conventions.

In the notation often associated with Dragt, products of Lie
transformations such as
\begin{equation}
    \e^{\lie{f_A}}\e^{\lie{f_B}}\vec z\,.
\end{equation}
are coordinate-substitution operators, so $A$ is encountered before $B$ in
the corresponding physical sequence.  If instead one writes maps with
their phase-space arguments, the same sequence is
\begin{equation}
    \mathcal{M}_B\bigl(\mathcal{M}_A(\vec z)\bigr)
    = \left(\mathcal{M}_B\circ\mathcal{M}_A\right)(\vec z)\,.
\end{equation}
There is no disagreement about composition: one expression acts on
functions (including the coordinate functions), while the other acts on
points.  Whenever an ordering is ambiguous, displaying either the argument
or the composition symbol $\circ$ removes it.
